\chapter{KI-Agenten}\label{ch:K05_Agenten}

\begin{myremark}
Während einfache LLMs \enquote{nur} auf Eingaben antworten, können \textbf{Agenten}
selbstständig handeln: Sie planen Aufgaben, nutzen Werkzeuge und führen mehrstufige
Prozesse durch. Dieses Kapitel erklärt, wie Agenten funktionieren und wo sie
eingesetzt werden.
\end{myremark}

% ─────────────────────────────────────────────────────────────────────────────
\section{Was ist ein KI-Agent?}

\begin{mydefinition}
Ein \textbf{KI-Agent} ist ein KI-System, das:
\begin{itemize}
    \item ein \textbf{Ziel} erhält (statt nur eine Frage),
    \item selbstständig \textbf{Teilschritte plant},
    \item \textbf{Werkzeuge nutzt} (Websuche, Code-Ausführung, Datenbankzugriff etc.),
    \item \textbf{iterativ vorgeht} und das Ergebnis jedes Schritts für den
          nächsten Schritt verwendet.
\end{itemize}
\end{mydefinition}

\begin{myexample}[– Schüler-Assistent]
Eine Studentin möchte einen Aufsatz über den Klimawandel schreiben. Sie gibt
einem Agenten folgenden Auftrag:

\begin{quote}
    \textit{Erstelle mir einen fundierten 5-seitigen Aufsatz über die Ursachen
    des Klimawandels. Nutze aktuelle wissenschaftliche Quellen und zitiere sie
    korrekt im APA-Format.}
\end{quote}

Der Agent führt folgende Schritte aus:
\begin{enumerate}
    \item Websuche nach aktuellen wissenschaftlichen Artikeln zum Klimawandel.
    \item Auswahl der relevantesten Quellen.
    \item Strukturierung des Aufsatzes (Gliederung erstellen).
    \item Verfassen der einzelnen Abschnitte.
    \item Einfügen der Zitate im APA-Format.
    \item Abschlusskontrolle und Überarbeitung.
\end{enumerate}
\end{myexample}

% ─────────────────────────────────────────────────────────────────────────────
\section{Wie funktionieren Agenten? – Das ReAct-Prinzip}

\begin{mydefinition}
Das \textbf{ReAct-Prinzip} (Reason + Act) beschreibt den Grundmechanismus
vieler KI-Agenten:

\begin{enumerate}
    \item \textbf{Reason (Denken)}: Der Agent analysiert die aktuelle Situation
          und überlegt, welcher Schritt als nächstes sinnvoll ist.
    \item \textbf{Act (Handeln)}: Der Agent führt eine Aktion aus (Websuche,
          Code schreiben, API aufrufen etc.).
    \item \textbf{Observe (Beobachten)}: Der Agent beobachtet das Ergebnis.
    \item Zurück zu Schritt 1: Der Zyklus wiederholt sich, bis das Ziel erreicht ist.
\end{enumerate}
\end{mydefinition}

\begin{myexercise}[– ReAct nachspielen]
\textbf{Spielen Sie selbst einen Agenten!} Sie erhalten folgende Aufgabe:

\begin{quote}
    \textit{Finde heraus, welches Tier in der Schweiz am häufigsten als
    Haustier gehalten wird, und erstelle eine kurze Tabelle mit den Top-3-Tieren.}
\end{quote}

Ohne das Internet zu nutzen:
\begin{enumerate}
    \item \textbf{Plan}: Welche Schritte würden Sie unternehmen?
    \item \textbf{Simulation}: Welche Werkzeuge würde ein echter Agent nutzen?
    \item \textbf{Ergebnis}: Erstellen Sie die Tabelle basierend auf Ihrem Wissen.
    \item \textbf{Reflexion}: Was hätte ein Agent besser gemacht?
\end{enumerate}
\end{myexercise}

% ─────────────────────────────────────────────────────────────────────────────
\section{Agentenframeworks und -tools}

\begin{myoverview}[Bekannte Agentenframeworks]
\begin{itemize}
    \item \textbf{OpenAI Operator}: Agent von OpenAI, der Browser steuern und
          Formulare ausfüllen kann.
    \item \textbf{Claude Computer Use}: Anthropics Agent, der den Computer
          direkt bedienen kann (Maus, Tastatur).
    \item \textbf{GitHub Copilot}: KI-Assistent für Programmierende, integriert
          in GitHub und VS Code. Kann Code schreiben, reviewen und Bugs finden.
    \item \textbf{LangChain}: Framework zum Bauen eigener Agenten (technisch).
    \item \textbf{AutoGPT}: Open-Source-Agent, der selbstständig Ziele verfolgt.
    \item \textbf{OpenClaw / Rent-a-Human}: Beispiele für KI-gestützte
          Plattformen, die menschliche Arbeit automatisieren oder koordinieren.
\end{itemize}
\end{myoverview}

\begin{myexercise}[– GitHub Copilot erleben]
Falls in Ihrer Schulumgebung GitHub Copilot verfügbar ist, erkunden Sie
folgende Funktionen:

\begin{enumerate}
    \item Öffnen Sie eine leere Datei in VS Code und schreiben Sie einen Kommentar:
          \texttt{\# Funktion, die die Fibonacci-Zahlen bis n berechnet}
          Beobachten Sie, was Copilot vorschlägt.
    \item Bitten Sie Copilot im Chat: \textit{Erkläre diesen Code Zeile für Zeile.}
    \item Bitten Sie Copilot, einen Bug in folgendem Code zu finden:
\begin{verbatim}
def addiere(a, b):
    return a - b
\end{verbatim}
\end{enumerate}

Diskutieren Sie: Verändert ein KI-Assistent wie Copilot, wie man Programmieren lernt?
Welche Fähigkeiten bleiben wichtig?
\end{myexercise}

% ─────────────────────────────────────────────────────────────────────────────
\section{NZZ-Artikel-Arena: Agenten und die Zukunft der Arbeit}

\begin{myattention}
Für diese Übung erhalten Sie NZZ-Artikel zu zwei kontroversen Themen aus dem Bereich
der KI-Agenten und der Plattformarbeit. Ziel ist eine strukturierte Debatte.
\end{myattention}

\begin{myprocedure}[Ablauf der NZZ-Artikel-Arena]
\begin{enumerate}
    \item \textbf{Vorbereitung (10 Min.)}: Jede Gruppe liest ihren Artikel und
          identifiziert die Hauptargumente.
    \item \textbf{Positionierung (5 Min.)}: Jede Gruppe formuliert eine klare
          Position (Pro / Contra) zum Thema.
    \item \textbf{Arena (15 Min.)}: Die Gruppen debattieren: Pro-Seite
          argumentiert, Contra-Seite widerspricht, dann Rollentausch.
    \item \textbf{Reflexion (10 Min.)}: Was sind die stärksten Argumente
          auf beiden Seiten? Gibt es einen Konsens?
\end{enumerate}
\end{myprocedure}

\subsection*{Arena-Thema 1: OpenClaw / KI-Agenten als Arbeitsersatz}

\begin{myexercise}[– Artikel-Arena: Agenten und Jobs]
Lesen Sie den folgenden NZZ-Artikel über OpenClaw und KI-Agenten als
Arbeitsersatz:

\begin{center}
    % TODO: PDF-Datei einfügen
    \includepdf[pages=-,scale=0.9]{Grundlagen_Info/14_KI_und_Prompting/Skript/Figures/NZZ_OpenClaw.pdf}
\end{center}

\textbf{These}: \textit{KI-Agenten sollten repetitive Büroarbeit vollständig
übernehmen, um Menschen für kreativere Aufgaben zu befreien.}

\begin{itemize}
    \item \textbf{Pro-Gruppe}: Sammeln Sie Argumente, die diese These unterstützen.
    \item \textbf{Contra-Gruppe}: Sammeln Sie Argumente, die diese These widerlegen.
\end{itemize}
\end{myexercise}

\subsection*{Arena-Thema 2: Rent-a-Human / Crowdworking}

\begin{myexercise}[– Artikel-Arena: Rent-a-Human]
Lesen Sie den folgenden NZZ-Artikel über Plattformen wie \enquote{Rent-a-Human}
und Crowdworking:

\begin{center}
    % TODO: PDF-Datei einfügen
    \includepdf[pages=-,scale=0.9]{Grundlagen_Info/14_KI_und_Prompting/Skript/Figures/NZZ_RentAHuman.pdf}
\end{center}

\textbf{These}: \textit{KI-Trainingsdaten werden oft von schlecht bezahlten
Arbeitnehmenden in Ländern des Globalen Südens produziert – das ist modernes
digitales Kolonialismus.}

\begin{itemize}
    \item \textbf{Pro-Gruppe}: Ist diese Kritik berechtigt? Welche Belege
          liefert der Artikel?
    \item \textbf{Contra-Gruppe}: Welche Gegenargumente gibt es?
          Ist Crowdworking per se schlecht?
\end{itemize}
\end{myexercise}

% ─────────────────────────────────────────────────────────────────────────────
\section{Grenzen und Risiken von Agenten}

\begin{mysummary}
Kritische Aspekte von KI-Agenten:
\begin{itemize}
    \item \textbf{Kontrollverlust}: Ein Agent, der selbstständig handelt,
          kann Fehler machen, die schwer rückgängig zu machen sind
          (z.\,B. unbeabsichtigt eine E-Mail versenden).
    \item \textbf{Sicherheit}: Agenten, die auf das Internet zugreifen,
          können durch Prompt-Injection-Angriffe manipuliert werden.
    \item \textbf{Autonomie vs. Aufsicht}: Wie viel Autonomie soll ein Agent haben?
          Wer ist verantwortlich, wenn ein Agent einen Fehler macht?
    \item \textbf{Datenschutz}: Agenten greifen oft auf persönliche Daten,
          Kalender, E-Mails etc. zu – das birgt Datenschutzrisiken.
\end{itemize}
\end{mysummary}

\begin{myexercise}[– Szenario: Agent macht Fehler]
Diskutieren Sie folgendes Szenario in der Gruppe:

\textit{Ein Unternehmen setzt einen KI-Agenten ein, der selbstständig
E-Mails im Namen der Chefin beantwortet. Der Agent missversteht eine Anfrage
und sagt einem wichtigen Kunden ab, obwohl die Chefin eigentlich zustimmen wollte.
Der Kunde bricht daraufhin den Vertrag ab.}

\begin{itemize}
    \item Wer trägt die Verantwortung: die Chefin, das Unternehmen oder
          der KI-Anbieter?
    \item Welche Kontrollmechanismen hätten diesen Fehler verhindert?
    \item Ab welchem Autonomiegrad sollten Agenten immer menschliche
          Bestätigung einholen?
\end{itemize}
\end{myexercise}

\begin{mychallenge}[Eigenen Mini-Agenten bauen]
\begin{myattention}
\textbf{Optional – Technische Challenge}: Diese Aufgabe erfordert Python-Kenntnisse
und einen OpenAI-API-Key (kostenpflichtig).
\end{myattention}

Erkunden Sie das Python-Framework \texttt{LangChain} (\url{https://langchain.com})
und bauen Sie einen einfachen Agenten, der:
\begin{enumerate}
    \item Eine Websuche durchführen kann.
    \item Das Ergebnis zusammenfasst.
    \item Die Antwort auf Deutsch ausgibt.
\end{enumerate}

Dokumentieren Sie Ihren Code und Ihre Beobachtungen.
\end{mychallenge}
